States, positivity and the adjoint are those of Premise \claimref{synthesis.premise-positivity-of-values}; the pairs are those of Proposition \claimref{tower.ladder-pairs-satisfy-car-vectors} and the tests those of Theorem \claimref{synthesis.set-and-clear-of-a-bit-generate-the-hyperbolic-plane}.

\emph{(i).} A state on one rung is fixed by $a = \omega(e)$, $c = \omega(f)$, $b = \omega(fe)$. Reality gives $c = 0$, since $f^{\dagger} = -f$. With $p = fe$, the elements $e$ and $p$ are self-adjoint, anticommute and square to $1$, so $u = ae + bp$ has $u^{2} = a^{2} + b^{2}$, and $\omega(x) = 2\langle\rho x\rangle_{0}$ with $\rho = \tfrac12(1 + u)$. If $a^{2} + b^{2} \le 1$ then $\rho$ is a nonnegative combination of the two self-adjoint idempotents $\tfrac12(1 \pm u/|u|)$ over the reals (or $\rho = \tfrac12$ when $u = 0$), so $\omega$ is positive; if $a^{2} + b^{2} > 1$ then $x = u - (a^{2} + b^{2})$ is self-adjoint with $\omega(xx^{\dagger}) = (a^{2} + b^{2})(1 - (a^{2} + b^{2})) < 0$. So positivity is $a^{2} + b^{2} \le 1$; with integer $a$, $b$ this leaves $(0,0)$, $(\pm 1, 0)$, $(0, \pm 1)$, five states, and $\omega(\mathrm{on}) = \tfrac12(1 + b)$. The script finds exactly these five by testing all $27$ integer triples in $\{-1, 0, 1\}^{3}$ against (P) (values outside that range violate $a^{2} + b^{2} \le 1$ or $c = 0$). The ladder's integer algebra is that of Theorem \claimref{synthesis.integer-clifford-algebras-contain-no-switch}. For the switch form, $\omega(vw) = \tfrac12(1 + b)$ and $\omega(v) = \tfrac12(a + c)$ must be integers, so $b = \pm 1$, and then $a = 0$: on and off. On $m$ rungs, the $P_{s}$ are products of the $v_{j}$, $w_{j}$, so an integer-valued state gives them nonnegative integers with sum $1$: one string $s$ has $\omega(P_{s}) = 1$. Then $\omega((1 - P_{s})(1 - P_{s})^{\dagger}) = 0$, and by the Cauchy and Schwarz inequality for the positive form $\omega(xy^{\dagger})$, $\omega(x) = \omega(P_{s}xP_{s})$ for all $x$; since $P_{s}xP_{s}$ is a multiple of $P_{s}$ ($P_{s}$ is a product of one test per rung, and each rung's algebra is the $2 \times 2$ matrices), $\omega$ is the state of the string. Its values on tests are $0$ or $1$ and it vanishes on each $e_{j}$.

\emph{(ii).} The even subalgebra of $\Cl_{\Q}(2,2)$ has dimension $8$ and is generated by the six bivectors; the script computes the space of even elements commuting with all six as the kernel of a linear map and finds dimension $2$; $1$ and $\Gamma$ lie in it. For the example the script verifies each stated property exactly: $yy^{\dagger}$ is an even idempotent; $X^{2} = 1$, $X^{\dagger} = X$, $X$ even; $UU^{\dagger} = 1$, $U$ even; the weights $\tfrac{9}{25}$, $\tfrac{16}{25}$ of the two record tests and their perfect correlation with the system's tests; $\Gamma$ acts as the same sign on both record components and the total number is $2$ with zero variance; $O$ is even and self-adjoint with value $\tfrac{24}{25} = 2 \cdot \tfrac35 \cdot \tfrac45$, and value $0$ in the state whose density is the sum of the two record-diagonal blocks. An even unitary commutes with $\Gamma$, so it cannot move a state from one sector to the other: the two outcomes of a record written by an even coupling on a state in one sector are in that sector. For $U'$: it is unitary since $e_{3}^{2} = 1$ and $e_{3}$ commutes with $\mathrm{on}_{1}$; its part $\mathrm{off}_{1}$ is even and its part $\mathrm{on}_{1}e_{3}$ is odd; the script exhibits an even state it sends to a state that is not even, the case of a record of an element that is not parity-homogeneous in Theorem \claimref{synthesis.records-keep-even-states-even-only-for-parity-homogeneous-observables}.

\emph{(iii).} With $h = v_{1}w_{2} - v_{2}w_{1}$ and $\Pi_{1} = P_{(1,0)} + P_{(0,1)}$, $h^{2} = -\Pi_{1}$ and $h^{\dagger} = -h$, so $R = (1 - \Pi_{1}) + \tfrac35\Pi_{1} + \tfrac45 h$ is even with $RR^{\dagger} = (1 - \Pi_{1}) + (\tfrac{9}{25} + \tfrac{16}{25})\Pi_{1} = 1$, and the script confirms $R\tilde{R} = 1$ and the weights $0, \tfrac{9}{25}, \tfrac{16}{25}, 0$ of $x \mapsto \omega_{(1,0)}(R^{\dagger}xR)$ on the four strings. A function from strings to strings carries the point weight on $(1,0)$ to a point weight; the weights after $R$ are not a point weight.

\emph{(iv).} By clause (ii) of Theorem \claimref{synthesis.state-of-a-new-rung-unique-symmetric-extension-is-the-half-weight}, the invariant extension vanishes on every blade containing a new generator; applied to two new rungs in turn it is $\omega \otimes$ (the normalised scalar part), and $P_{ab}$ has scalar part $\tfrac14$. Clause (iii) of that theorem is the statement about labels.

The Reading is argued from (i) to (iv) and the stored statements it cites; it is not a theorem.
